\section{Exact Minimum-Investment Solvers}
\label{sec:method}

For every rule and cap we compute the exact optimum, not a heuristic. For a score rule $f$ let $\tau=\max_{j\ne i}f(X_j)$ and $b=\tau+\varepsilon$, so that first place means $f(y)\ge b$; let $H_t=c_t-x_t$ be the headroom on task $t$. $L_1$, $K$ and $L_\infty$ are three separate optimisations, so the $K$-minimal and the $L_1$-minimal raise generally differ. All rules are monotone in the target's own scores, hence first place is reachable under a cap if and only if the target raised to the cap $c$ is first; every solver uses this as its feasibility test.

\begin{proposition}[Mean]\label{prop:mean}
If $\sum_tH_t\ge nb-\sum_tx_t$ then $L_1^\ast=\max(0,nb-\sum_tx_t)$ and $K^\ast$ is the smallest $k$ such that the $k$ largest headrooms sum to at least $L_1^\ast$; otherwise first place is infeasible. If no cap binds, $L_\infty^\ast=b-\bar x$ and $L_1^\ast=n\,L_\infty^\ast$.
\end{proposition}
Under the mean, spreading the investment costs exactly as many points as concentrating it: targeting buys a smaller $K$, not fewer points.

\begin{proposition}[Minimum]\label{prop:min}
First place is feasible iff $c_t\ge b$ for all $t$. Then $L_1^\ast=\sum_t(b-x_t)_+$, $K^\ast=|\{t:x_t<b\}|$, $L_\infty^\ast=\max_t(b-x_t)_+$, so $L_1^\ast\le nL_\infty^\ast$.
\end{proposition}

\begin{proposition}[Median, $n$ odd]\label{prop:median}
With $m=(n+1)/2$, first place is feasible iff at least $m$ tasks have $c_t\ge b$. Then $K^\ast=\max\bigl(0,m-|\{t:x_t\ge b\}|\bigr)\le m$ and $L_1^\ast$ is the sum of the $m$ smallest shortfalls $(b-x_t)_+$ among tasks with $c_t\ge b$.
\end{proposition}
Proofs are in Appendix~\ref{app:proofs}. For even $n$ (Open LLM Leaderboard, $n=6$) the median is the mean of the two middle scores and the solver enumerates the pair of threshold levels. Power means with exponent $p<1$ (geometric: $p\to0$; harmonic: $p=-1$, both on $y+1$) are $\frac1n\sum_tg(y_t)\ge g(b)$ with $g$ increasing and concave; $L_1^\ast$ is then a water-filling problem (bisection on the common level $y_t+1$, clipped to $[x_t,c_t]$) and $K^\ast$ is greedy in $g(c_t)-g(x_t)$. The $t$-trimmed mean has no closed form and is solved as a mixed-integer program (MILP).

\paragraph{Rank rules.}
A raise of the target on task $t$ to a level just above a rival's score passes every model at or below that score, so only the levels $\{s_{kt}\}$ (tie, half credit) and $\{s_{kt}+\varepsilon\}$ (strict pass) matter, at most one per task. Borda, plurality and the Condorcet-winner requirement are then exact 0/1 programs, solved with HiGHS (SciPy) with lazily generated rival constraints. The Condorcet winner (the target wins strictly more than $n/2$ tasks against every rival) is sufficient for first place under any Condorcet-consistent rule, so its minimum investment is an \emph{upper bound} for that family; the reported ``rule rank'' for it is a Copeland proxy.

\paragraph{Verification.}
Propositions~\ref{prop:mean}--\ref{prop:median} are proved and also compared with the solvers on all 79 HELM targets (Section~\ref{sec:exp-theory}). The trimmed-mean and rank-rule MILPs are exact by construction but come with no separate proof; they are checked against independent formulations (an assignment formulation of the trimmed mean, a grid search for the even median) and against brute-force enumeration on small random instances.
